% 第8章 竞争相互作用与低维性（Competing interactions and low dimensionality）
\chapter{竞争相互作用与低维性}

本章考察竞争相互作用与低维性使固体中出现某些极为精微、复杂、有时甚至有用的磁
行为的一些方式。竞争相互作用与低维性都天然地出现在某些体系中：例如有些晶体生长
的方式使磁矩在链中强耦合、或占据在阻挫晶格的格点上。然而这些特征也可以人为引入，
例如制备铁磁多层膜，或用电子束光刻制造铁磁线。其中许多专题目前正被世界各地的研究
小组积极研究。本章叙述这些进展中的一部分——其中许多仍笼罩着尚未揭开的谜。由于
很多专题既新又复杂，有些讨论难免过于简化；有兴趣的读者可查阅章末延伸阅读中的
文献以了解详情。尽管如此，我希望传达出若干问题的风味——在写作之时，它们正给
凝聚态物理学家带来巨大的挑战。

我们从阻挫（frustration）开始——它可以导致包括自旋玻璃在内的多种基态——然后
继续讨论低维性的效应。

\section{阻挫}

在某些晶格中，不可能同时满足体系中的所有相互作用以找到基态。这常常导致没有唯一
的基态，而是存在许多相似的低能态，体系的不满足感（指能量未被极小化）被尽可能地
分摊。此时称体系表现阻挫。例如考虑图 8.1 的晶格，其中只有最近邻反铁磁相互作用
起作用。在正方晶格（图 8.1(a)）上，让最近邻自旋反平行的要求很容易满足；但在三角
晶格上事情就不那么简单了：如图所示，若两个相邻自旋反平行排列，第三个自旋便陷入
两难——无论怎么选，两个邻居中总有一个的能量未被极小。体系因此无法达到完全满足其
微观约束的状态，但拥有大量同样不满足的态。结果，阻挫体系表现出亚稳性、磁滞效应
（依赖样品的磁历史或热历史）以及趋向平衡的时间弛豫。

\begin{figure}
\centering
\includegraphics[width=0.45\textwidth]{fig8_01.pdf}
\caption{（a）正方晶格与（b）三角晶格上的反铁磁最近邻相互作用。三角晶格表现阻挫：
第三个格点上的自旋无论怎样取向，都无法与另两个自旋同时满足反铁磁最近邻相互作用的
要求。}
\end{figure}

于是某些体系中晶格的几何可以阻挫自旋的有序化。二维中，三角晶格与共角 kagomé
晶格（见图 8.2）上的海森伯自旋表现这一效应；三维中被研究得最充分的体系具有
烧绿石（pyrochlore）结构——磁性离子占据共角四面体构成的晶格。据信这些体系不
有序，而是表现具有宏观简并的经典基态——有时称为协同顺磁性（cooperative
paramagnetism）：一切温度下自旋之间都只有短程关联。这些体系中的一些还被认为
具有无色散的自旋波支——称为零模（zero mode），它强烈影响低温热力学行为，导致
自旋涨落一直持续到零温。

\begin{figure}
\centering
\includegraphics[width=0.4\textwidth]{fig8_02.pdf}
\caption{kagomé 晶格的一部分。}
\end{figure}

\section{自旋玻璃}

自旋玻璃（spin glasses）已在 5.5 节介绍。该节把自旋玻璃定义为：随机、混合相互
作用的磁性体系，其特征是在明确的温度 $T_{\mathrm{f}}$ 处自旋发生随机然而协同的
冻结，冻结温度以下出现高度不可逆的亚稳冻结态，而没有通常的磁长程序。现在我们
更细致地考察这一定义的各个部分，先看可以用几种方式引入的随机性。

其一是位置随机性（site-randomness），可在合金中实现。常被研究的自旋玻璃是
$Cu_{1-x}Mn_x$（$x\ll1$）：少量 Mn 完全随机地替代进入 Cu 基体、没有任何短程
有序。这直接导致非磁 Cu 基体中磁性 Mn 离子间距离随机（见图 5.1(c)）。取金属间
化合物如 $GdAl_2$ 使其非晶化（熔化后从熔体极速冷却，晶格被摧毁）也可以人工设计
磁离子间的随机距离。

另一种可能是键随机性（bond-randomness）：最近邻相互作用在 $+\mathsf{J}$ 与
$-\mathsf{J}$ 之间变化。可以把磁性离子固定在规则的晶格阵列上，同时调制磁性离子
之间的间接交换相互作用来实现。这在 $Rb_2Cu_{1-x}Co_xF_4$ 中实现：Cu 与 Co 的
有效自旋都是 $\frac{1}{2}$（Co 上的晶体场把 $S=\frac{3}{2}$ 能级分裂，使低温下
只有最低的 $S=\frac{1}{2}$ 二重态被占据）。晶体场还使自旋平行或反平行于 z 轴
（单轴单离子各向异性，见 3.2.2 节）——称其具有伊辛性格。然而磁性离子间超交换
相互作用的大小与符号关键取决于耦合是 Co--Co、Co--Cu 还是 Cu--Cu，以及 Cu 上
占据的是哪个轨道。净结果是离子间的键可以取不同的交换 $\mathsf{J}$ 值，这些键
随机分布在整个样品中。

自旋玻璃固有的随机性很重要，竞争相互作用的存在同样重要。随机位置自旋玻璃中
磁矩间距离的分布导致竞争相互作用，因为相互作用是 RKKY 型的，其符号（铁磁或
反铁磁）依赖自旋间距。另一要素是磁各向异性——来自单离子各向异性或
Dzyaloshinsky--Moriya 相互作用。非晶材料中这些各向异性逐点不同，因此所谓的
非晶磁体常具有随机各向异性：磁化强度的“易轴”可以局域变化。

随机键自旋玻璃中竞争相互作用自动存在——不同的键以不同方式“拉扯”系统。这些
竞争相互作用导致阻挫：像 kagomé 晶格那样，存在多重简并的基态。自旋玻璃共享这种
多重简并基态，但还表现一种新效应——协同冻结转变，现在就来描述。

高温下一切磁性系统的行为由热涨落主导，自旋玻璃中所有自旋彼此独立。自旋玻璃从
高温冷却时，独立自旋减速并结成局域关联的单元——称为团簇（clusters）。不在团簇
里的自旋参与团簇之间的相互作用。温度降到 $T_{\mathrm{f}}$ 时，团簇的涨落逐渐
减慢。自旋间的相互作用变得更加长程，每个自旋越来越“感知”到其周围逐渐扩大的
区域内的自旋。在 $T_{\mathrm{f}}$ 处系统找到它的众多基态之一并冻结。这一过程
尚未被完全理解，但看起来是一种协同相变。然而这不是向磁有序态的相变：若体系有
长程磁有序，散射实验中应出现磁布拉格峰——实测没有。$T_{\mathrm{f}}$ 以下基态
显得“玻璃态”：具有亚稳性与缓慢的弛豫行为。$T_{\mathrm{f}}$ 以下场冷却与零场
冷却磁化率出现分岔，反映这种亚稳性（见图 8.3）。

自旋玻璃行为的一个标志是交流磁化率实部 $\chi(\omega)$ 在 $T_{\mathrm{f}}$ 附近
的尖峰。该技术用频率 $\omega$ 的很小的交变磁场测量磁化率，有时同时加恒定（直流）
磁场。峰位随 $\omega$ 略有移动。与吸收相关的 $\chi(\omega)$ 虚部在 $T_{\mathrm{f}}$
附近突然出现。冻结过程的涨落动力学可用交流磁化率、中子自旋回波（在相对长时间
尺度上测量时间关联函数的方法）与 $\mu$SR 研究。这些技术表明：$T_{\mathrm{f}}$
以上形成的团簇具有很宽的弛豫时间分布。样品冷却时，团簇长大并以依赖各自尺寸与
性质的速率涨落。团簇之外的自旋快速涨落，但它一加入正在生长的团簇（有点像“恶意
收购”！）就被迫减速、与其新东家同步。在 $T_{\mathrm{f}}$ 处随机各向异性变得
重要，系统冻结成团簇固定的随机取向。$T_{\mathrm{f}}$ 以下也观察到很宽的弛豫时间
范围，表明仍有一些自由自旋或小的超顺磁（见下文）团簇存在。自旋玻璃中的协同自旋
冻结仍未被完全理解。

\begin{figure}
\centering
\includegraphics[width=0.5\textwidth]{fig8_03.pdf}
\caption{$Cu_{1-x}Mn_x$（x=0.0108 与 x=0.0202）的静态磁化率。零场冷却后，在
0.59 mT 的小场中升温测量 $\chi$（(b) 与 (d)）。若样品在 0.59 mT 中场冷却，$\chi$
可逆（(a) 与 (c)）。改编自 Nagata 等 1979。}
\end{figure}

$Cu_{1-x}Mn_x$ 自旋玻璃通常在 x 为百分之几以下时研究。浓自旋玻璃中行为大不相同：
磁最近邻数目更多（可因化学偏聚而增强），导致出现大的、有序的磁团簇并主导磁行为。
系统于是称为团簇玻璃（cluster glass，有时称 mictomagnet）——具有自旋玻璃的某些
特征，但长程磁有序始终潜伏在背后、伺机而发（术语谓之长程磁有序是“萌芽态”
（incipient））。有时高温态实际上可以是有序铁磁体而非顺磁体，系统实际冻结到低温
非磁无序态——这称为重入自旋玻璃（re-entrant spin glass），其行为据信与温度依赖
的随机各向异性或与冻结相联系的某种各向异性有关。磁性离子浓度很高时，系统逼近
渗流极限——可以找到一条由最近邻连键构成的、贯穿整个样品的路径，路径上每个离子
都是磁性的；这使样品具有长程磁有序。这些不同态之间的转变可见于 $Au_{1-x}Fe_x$
的相图（图 8.4）。

\begin{figure}
\centering
\includegraphics[width=0.5\textwidth]{fig8_04.pdf}
\caption{$\mathrm{Au}_{1-x}\mathrm{Fe}_x$ 的相图。P=顺磁，SP=超顺磁，SG=自旋
玻璃，CG=团簇玻璃，F=铁磁。取自 Coles 等 1978。}
\end{figure}

\section{超顺磁性}

若铁磁颗粒足够小，它们将是单畴的（见 6.7.7 节），因为形成畴壁的能量代价超过
退磁能的节省。小的单畴铁磁颗粒的磁化强度常被约束在平行或反平行于特定方向。原因
可以是磁晶各向异性、形状各向异性（与退磁能及颗粒形状相联系），或其他多种因素。
下面假定颗粒的能量密度含 $K\sin^{2}\theta$ 项——$\theta$ 是磁化强度与该特定
方向的夹角，K 是量化这一各向异性能量密度的常数。于是 $\theta=0$ 或 $\pi$ 时
能量最低（见图 8.5）。体积 V 的颗粒把磁化强度从 $\theta=0$ 翻转到 $\pi$（或
反之）需要激活能 $\Delta E=KV$。对非常小的颗粒，$KV$ 与 $k_{\mathrm{B}}T$ 相比
很小，热涨落便能轻易完成这种翻转。

\begin{figure}
\centering
\includegraphics[width=0.35\textwidth]{fig8_05.pdf}
\caption{磁性颗粒的能量密度含 $K\sin^{2}\theta$ 项；$\theta=0$ 或 $\pi$ 时能量
最低。}
\end{figure}

考虑非磁基体中这些小铁磁颗粒的分布，并设颗粒相距足够远、颗粒间相互作用可忽略。
$k_{\mathrm{B}}T\gg KV$ 时系统表现得像顺磁体——只不过独立“矩”不是原子矩，而是
铁磁颗粒内的一大组矩。这样的系统称为超顺磁体（superparamagnet）。高温时颗粒上
的矩能快速涨落。颗粒上矩的弛豫时间 $\tau$ 为
\begin{equation*}
\tau = \tau_0\exp\left(\frac{KV}{k_{\mathrm{B}}T}\right),\tag{8.1}
\end{equation*}
$\tau_0$ 典型为 $10^{-9}$ s。样品冷却时涨落减慢（$\tau$ 增大，见图 8.6）；当
$\tau$ 远长于所用实验技术的测量时间 t 时，系统显得静态。若定义“远长于”为
$\tau>\alpha t$（$\alpha=100$），则低于阻塞温度（blocking temperature）
$T_{\mathrm{B}}=KV/k_{\mathrm{B}}\ln(\alpha t/\tau_0)$ 时，每个磁性颗粒看起来被
锁进它的两个极小之一。

\begin{figure}
\centering
\includegraphics[width=0.42\textwidth]{fig8_06.pdf}
\caption{按式 8.1，弛豫时间 $\tau$ 随温度 T（以 $k_{\mathrm{B}}/KV$ 标度）的依赖。
温度降低时涨落减慢（$\tau$ 增大）。}
\end{figure}

实验测量时间 t：非弹性中子散射为 $10^{-12}$--$10^{-10}$ s；穆斯堡尔为
$10^{-10}$--$10^{-7}$ s（穆斯堡尔跃迁的衰变过程 $\sim10^{-7}$ s）；$\mu$SR 为
$10^{-10}$--$10^{-5}$ s（可测比例的 $\mu$ 子能活到 $\sim10\tau_{\mu}$，
$\tau_{\mu}=2.2\ \mu$s 是 $\mu$ 子平均寿命）；交流磁化率典型探测
$10^{-1}$--$10^{-5}$ s。由于对 $\alpha t/\tau_0$ 是对数依赖，把 $\alpha$、t 或
$\tau_0$ 改变几个量级，$T_{\mathrm{B}}$ 也只有相对较小的变化。若超顺磁体系中
颗粒有一系列尺寸，它们将在不同温度“阻塞”。

超顺磁体与自旋玻璃有某些相似，但多方面截然不同：超顺磁体中磁性颗粒之间的相互
作用不重要，自旋玻璃中却至关重要；自旋玻璃表现协同相变，超顺磁体则表现超顺磁
颗粒的逐渐阻塞。超顺磁性在技术上很重要，因为许多重要记录材料是颗粒型的。这一点
至关重要：磁带被期待存储信息数年而非 $\mu$s 量级——这给颗粒的最小尺寸设了限。
研究岩石磁性时，有时需要考虑磁性在数百万年上的指数衰减。因此颗粒对超顺磁涨落的
稳定性依赖你所用的时标：同一颗粒在穆斯堡尔实验中稳定、在交流磁化率中却可能涨落；
可以稳定一个世纪，却在地质时标上衰减。

\section{一维磁体}

一维中 T>0 不可能有长程序，人们或许以为一维磁体相当乏味无趣。事实恰恰相反！
一维性意味着可能存在复杂的激发——至今远未被完全理解。

\subsection{自旋链}

一维（d=1）的自旋线称为自旋链（spin chain）。单个自旋可以被约束在平行或反平行
于特定方向（伊辛自旋），或自由指向固定平面内任意方向（XY 自旋），或自由指向
任意方向（海森伯自旋）。若晶体结构使链之间保持足够距离，自旋链可以在晶体中近似
实现。晶体场导致的单离子各向异性可使磁矩表现得像伊辛自旋（D=1）、XY 自旋
（D=2）、海森伯自旋（D=3）或介于其间。一个常被研究的体系家族基于 ABX$_3$ 型
晶体：A 是一价非磁阳离子，B 是二价磁性阳离子，X 是卤素阴离子。这给出简单六角
晶格，过渡金属离子沿 c 方向构成链。例如 CsCoCl$_3$ 中各向异性把自旋约束在特定
方向，几乎表现为一维伊辛自旋链；KCuF$_3$ 表现为一维海森伯自旋链，许多带有机
配体的 Cu 盐亦然。

这些体系常在极低温下表现三维长程序——总存在某种小的链间相互作用把链耦合起来。
正因为这种链间耦合，CsCoCl$_3$ 在 21 K 以下表现三维长程磁有序。不过在向三维
行为过渡的温度之上，仍有很宽的温度区域，其磁行为是一维的。

链上每个自旋的自旋量子数取决于原子：$\mathrm{Cu^{2+}}$（3d$^9$）链为
$S=\frac{1}{2}$，$\mathrm{Mn^{2+}}$（3d$^5$）为 $S=\frac{5}{2}$；CsCoCl$_3$ 中
$\mathrm{Co^{2+}}$（3d$^7$）的基态具有有效自旋 $S=\frac{1}{2}$（$S=\frac{3}{2}$
自由离子基态被晶体场分裂，留下基态二重态与激发态）。下一节只考虑每个格点
$S=\frac{1}{2}$ 的链。

\subsection{自旋子}

使这些链有趣的不它们的有序（因为链间相互作用不够强时它们不表现有序），而是其
激发。如 6.6.1 节所示，三维海森伯磁体的激发是磁振子——玻色子（金属磁性材料中
还可能有 Stoner 激发，见 7.8 节）。每个磁振子是 S=1 的激发，能在非弹性散射实验
中与中子作用。

伊辛自旋链的激发与畴壁的产生相联系。不存在无能隙激发（即长波极限下能量趋于零
的激发；用粒子物理的行话叫“无质量戈德斯通模”），因为哪怕产生一条畴壁也要付出
有限能量（实际上激发必须成对地产生畴壁）。激发一旦产生便可沿链自由移动。激发
能量不依赖波矢，但若链不是完全伊辛型的（真实体系常如此），色散关系会有一些调制。

海森伯自旋链的激发称为自旋子（spinons）。它们具有自旋 $\frac{1}{2}$
（与自旋 1 的磁振子相对照），是费米子。色散关系为
\begin{equation*}
\hbar\omega = \pi|\mathsf{J}\sin(qa)|,\tag{8.2}
\end{equation*}
$\mathsf{J}$ 是反铁磁交换耦合，a 与 q 是沿链方向的晶格常数与波矢。式 8.2 即
图 8.7 中的粗线。可与式 6.58 的常规自旋波色散（取 $S=\frac{1}{2}$）相比——多了
一个因子 $\pi/2$。两种情形的激发都是无能隙的：$q\to0$（长波极限）时
$\omega\to0$。中子散射实验涉及自旋改变 1：在三维系统中这意味着单个磁振子的产生
或湮灭，在一维系统中则意味着两个自旋子的产生或湮灭。因此中子实验测量的是与一对
自旋子的产生或湮灭相联系的动量 $q=q_1+q_2$ 与能量
$\hbar\omega=\hbar\omega_1+\hbar\omega_2$，实验数据在式 8.2 与
$\hbar\omega=2\pi|\mathsf{J}\sin(qa/2)|$ 之间显示连续的激发（见图 8.7）。中子
散射实验已证实在某些一维反铁磁链中这些激发确实存在（见图 8.8）。（式 8.2 的推导
见 des Cloizeaux 与 Pearson 1962。）

\begin{figure}
\centering
\includegraphics[width=0.5\textwidth]{fig8_07.pdf}
\caption{一维反铁磁海森伯自旋链中自旋子激发的色散关系（粗线）。阴影区为中子散射
实验测得的激发连续谱。}
\end{figure}

\begin{figure}
\centering
\includegraphics[width=0.5\textwidth]{fig8_08.pdf}
\caption{（a）KCuF$_3$ 中自旋子的色散关系。实验采用飞行时间技术，直线为入射能量
148.9 meV、入射动量偏离 KCuF$_3$ 的 c$^*$ 方向 8° 的中子的散射轨迹；当轨迹与
连续态相交时发生散射。（b）20 K 下观察到的散射。非磁背景以虚线标出。取自
Tennant 等 1993。}
\end{figure}

\subsection{Haldane 链}

上一节指出，自旋 $\frac{1}{2}$ 反铁磁海森伯链的激发据信是无能隙的自旋子。这一
结果据信对半整数自旋链（$S=\frac{1}{2},\frac{3}{2},\frac{5}{2},\ldots$ 的链）
同样成立。Haldane 猜想：整数自旋链（$S=1,2,3,\ldots$）会有所不同——激发谱中
会出现能隙，其起因是基态中的非线性量子涨落。整数自旋的一维链因此称为 Haldane
链（Haldane chain），激发谱中的能隙称为 Haldane 能隙。半整数与整数自旋链的这一
根本差别与费米子和玻色子在交换下的差别相关；这种不同的交换对称性具有拓扑起源，
对激发的性质有戏剧性的影响。

Haldane 猜想的多数检验在 $Ni^{2+}$（S=1）离子链材料上进行，包括 CsNiCl$_3$、
$\mathrm{Ni(C_2H_8N_2)_2NO_2ClO_4}$ 与 $Y_2BaNiO_5$。这些 S=1 自旋链的激发谱
都如预言地拥有能隙。半整数反铁磁链是无能隙的——除非磁弹耦合打开自旋 Peierls
能隙，见下一节。

\subsection{自旋 Peierls 转变}

尽管自旋 $\frac{1}{2}$ 反铁磁链无能隙，它容易遭受一种与折磨一维金属的不稳定性
（见 7.9 节）类似的、能打开能隙的不稳定性。这发生在自旋 Peierls 转变（spin-Peierls
transition）。

\begin{figure}
\centering
\includegraphics[width=0.42\textwidth]{fig8_09.pdf}
\caption{（a）均匀海森伯反铁磁链与（b）交替链（基态是 q=0 的单态，单胞加倍）中
元激发的示意。改编自 Bray 等 1983。}
\end{figure}

这种本征晶格不稳定性的驱动力是一维电子结构与三维晶格振动（声子）之间的磁弹耦合。
耦合产生的原因是链的交换能量是相邻格点间距的函数。晶格畸变影响磁能量（见图
8.9）。“自旋 Peierls”之名反映了它与 Peierls 畸变（7.9 节讨论过）的相似性。

转变温度 $T_{\mathrm{SP}}$ 以上，每条链中有均匀的反铁磁近邻交换；$T_{\mathrm{SP}}$
以下出现弹性畸变，导致二聚化，进而产生两个不相等的交替交换常数。二聚化随温度降低
渐进增大，在零温达到极大。交替链在单态基态与最低三重激发带之间有能隙。能隙大小
与二聚化程度（即晶格畸变程度）相关：均匀链（零二聚化）时为零——回到无能隙自旋子
情形。于是磁化率 $\chi(T)$ 在 $T_{\mathrm{SP}}$ 处呈现“膝”，$T_{\mathrm{SP}}$
以下 $\chi$ 相当突然地跌落，对应能隙的打开（见图 8.10）。通常的 Peierls 畸变
（自旋 Peierls 转变的电子类比，见 7.9 节）发生在 $k_{\mathrm{B}}T_{\mathrm{P}}\sim E_{\mathrm{F}}\exp(-1/\lambda_{\mathrm{el-ph}})$
量级的温度（$\lambda_{\mathrm{el-ph}}$ 为电子--声子耦合常数），而自旋 Peierls
转变发生在 $k_{\mathrm{B}}T_{\mathrm{SP}}\sim|\mathsf{J}|\exp(-1/\lambda_{\mathrm{sp-ph}})$
（$\mathsf{J}$ 为相邻自旋间的交换相互作用，$\lambda_{\mathrm{sp-ph}}$ 为自旋--
声子耦合常数）。由于 $\mathsf{J}\ll E_{\mathrm{F}}$（例如 $\mathsf{J}$ 典型为 50 K，
$E_{\mathrm{F}}$ 典型为 500--5000 K），$T_{\mathrm{SP}}$ 总是远小于 $T_{\mathrm{P}}$。

表现自旋 Peierls 转变的材料非常少，因为反铁磁链常因链间耦合在低温下变为三维有序。
只有极少数材料，自旋--声子耦合能压倒链间自旋--自旋耦合、允许自旋 Peierls 基态
形成。例子包括 CuGeO$_3$（$T_{\mathrm{SP}}=14$ K）与若干有机体系，如
MEM(TCNQ)$_2$（$T_{\mathrm{SP}}=18$ K）与 TTF-CuS$_4$C$_4$(CF$_3$)$_4$
（$T_{\mathrm{SP}}=12$ K）。\fn{1}{MEM、TCNQ 与 TTF 是具有冗长化学名称的有机分子。}

\begin{figure}
\centering
\includegraphics[width=0.5\textwidth]{fig8_10.pdf}
\caption{有机自旋 Peierls 材料 MEM(TCNQ)$_2$ 的摩尔磁化率——它由有机分子 MEM 的
堆叠与有机分子 TCNQ 的堆叠构成。高温下磁化率符合均匀海森伯反铁磁链的模型（点线）；
降温时磁化率在自旋 Peierls 转变处随激发谱能隙的打开而迅速下降。极低温（T 降低时）
的回升来自缺陷的类居里（$\sim T^{-1}$）磁化率。取自 Lovett 等 2000。}
\end{figure}

\subsection{自旋梯}

在转向二维磁体之前，可以考虑一个介于一维磁体与二维磁体之间的体系。考虑两条平行的
自旋链，链间有键，链间耦合与链内耦合强度相当。这样的系统称为两腿自旋梯（two-leg
spin ladder，见图 8.11(a)）。也可以有三腿（图 8.11(b)）、四腿乃至 n 腿自旋梯。
有前景的实验体系包括 SrCu$_2$O$_3$ 与 La$_{1-x}$Sr$_x$CuO$_{2.5}$（都是两腿
自旋 $\frac{1}{2}$），以及 Sr$_2$Cu$_3$O$_5$（三腿自旋 $\frac{1}{2}$）。事实上
存在一个通式为 Sr$_{n-1}$Cu$_{n+1}$O$_{2n}$ 的体系（n 为奇数），它是 (n+1)/2 腿自旋
$\frac{1}{2}$ 体系，因为其结构含有沿宽度方向排布 (n+1)/2 个 Cu$^{2+}$ 离子的
CuO$_2$ 正方晶格条带。

\begin{figure}
\centering
\includegraphics[width=0.3\textwidth]{fig8_11.pdf}
\caption{（a）两腿梯；（b）三腿梯。（c）两腿梯的基态由梯子每一横档上的自旋单态
组成，因此两腿梯的激发谱有能隙。（d）往梯上掺入空穴会破坏单态，（e）但若空穴
配对，这一能量代价可以被最小化。}
\end{figure}

已知自旋 $\frac{1}{2}$ 两腿梯的激发谱有有限能隙——在“强横档”极限（横档耦合
$J_{\perp}$ 大于沿腿的耦合 $\mathsf{J}$）下一眼可见：基态就是每个横档上的自旋
单态（见图 8.11(c)）。要产生激发，必须把横档单态提升为横档三重态（耗能
$J_{\perp}$），故有能隙。$J_{\perp}=0$ 时系统是两条孤立的自旋 $\frac{1}{2}$ 链，
激发谱无能隙。然而据信只要 $J_{\perp}$ 非零——无论多小——能隙就出现。对 n 腿
梯，n 为偶数时情形相同；n 为奇数时，给定横档上的自旋配成单态后总剩下一个落单。
大 $J_{\perp}$ 时系统可映射为无能隙的自旋 $\frac{1}{2}$ 链。所以：偶数腿梯的
激发谱有能隙，奇数腿梯无能隙。中子、$\mu$SR 与输运实验的结果似乎支持这一点。

往自旋梯中掺杂空穴会破坏自旋梯中的单态（见图 8.11(d)）。对两腿梯，空穴配对在
能量上有利——它们可以共享同一横档，把“受损”单态从两个减少到一个
（图 8.11(e)）。因此空穴在两腿梯上配对是有利的——这展示了在两腿自旋梯中设计
出超导电性的可能。超导电性（$T_{\mathrm{c}}\sim14$ K）实际上已在
Sr$_{14-x}$Ca$_x$Cu$_{24}$O$_{21}$（有时称 [14-24-41] 或“电话号码”化合物）
中发现（x=13.6、5 GPa）。对自旋梯的兴趣在于：它们可以有激发谱能隙、可以变为
超导，却又是简单而定义明确的体系——喜欢在一维上工作的理论家可以尝试建模。因此
这些磁性体系可能照亮高 $T_{\mathrm{c}}$ 超导电性问题（见下面的 8.5 节）。

\section{二维磁体}

二维磁性常在化学式 $A_2BX_4$ 的体系中研究：与前面一样，A 是一价非磁阳离子，
B 是二价磁性阳离子，X 是卤素阴离子。晶体结构为四角，磁性离子位于二维正方晶格上。
典型材料是 $K_2NiF_4$——这些体系常被称为具有 $K_2NiF_4$ 结构。Mermin--Wagner--%
Berezinskii 定理表明：$d\leq2$ 时，各向同性体系中热涨落禁止非零温度下长程磁
有序的存在。但它对 T=0 基态不置一词。一维中，事实证明自旋 $\frac{1}{2}$ 海森伯
反铁磁体即使在 T=0 也没有长程序。二维海森伯反铁磁体的情形不那么明确，目前有
证据表明它在 T=0 确实表现长程磁有序。这一体系之所以极受关注，是因为二维海森伯
模型看来对理解高 $T_{\mathrm{c}}$ 铜氧化物超导体很重要。这些体系包含铜氧面
（见图 8.12），铜氧面对超导性质似乎至关重要。

\begin{figure}
\centering
\includegraphics[width=0.4\textwidth]{fig8_12.pdf}
\caption{$\mathrm{La_2CuO_4}$ 中的铜氧面。}
\end{figure}

例如，纯 $La_2CuO_4$ 是反铁磁绝缘体、具有 $K_2NiF_4$ 结构。$Cu^{2+}$ 离子
（3d$^9$）在 d 带有一个空穴、自旋 $\frac{1}{2}$，位于正方晶格上；每个
$Cu^{2+}$ 离子经氧离子与邻居隔开，氧传递 $Cu^{2+}$ 自旋之间的反铁磁超交换相互
作用（见图 8.12）。每格点一个电子，你会预期系统是金属，但空穴因关联而局域化。
用额外的空穴掺杂（以 $Sr^{2+}$ 替换部分 $La^{3+}$）导致最佳掺杂样品（约 20\%
的 $La^{3+}$ 被 $Sr^{2+}$ 替换）在 40 K 左右出现高温超导。可以用 $Ni^{2+}$
（S=1）替换 $Cu^{2+}$ 制备同结构材料，但它们即使掺杂也不超导——因此 $Cu^{2+}$
有某种特殊之处。特别令人惊讶的是这些磁性离子竟有助于促成超导电性——磁性通常会
摧毁超导电性：磁性离子产生的局域磁场充当对破坏者（pair-breaker），拆散承载
超流的库珀对。因此，二维 $S=\frac{1}{2}$ 正方晶格海森伯反铁磁体的磁性质可能与
高 $T_{\mathrm{c}}$ 超导电性问题相关——尽管这一论断未被证明：写作之时，尚无
公认的高 $T_{\mathrm{c}}$ 超导理论。

二维 $S=\frac{1}{2}$ 正方晶格海森伯反铁磁体因此备受关注，吸引了大量理论与实验
研究。温度降低时，系统并不有序，但短程有序的关联区域尺寸增大。这一尺寸称为自旋--
自旋关联长度 $\xi$，低温下指数发散，导致 T=0 的真正长程序。通过把二维 $S=\frac{1}{2}$ 正方晶格海森伯反铁磁体映射到一个称为二维量子非线性
sigma 模型的连续模型，可以计算 $\xi$ 随温度的行为。这一步骤超出本书范围，但为这一方法提供了大量支持的
实验结果示于图 8.13。结果由 $Sr_2CuO_2Cl_2$ 的非弹性中子散射得到——它是理想
二维 $S=\frac{1}{2}$ 正方晶格海森伯反铁磁体的极好近似：面间距比 $La_2CuO_4$
大，且直到低温保持四角。面间相互作用导致 $T_{\mathrm{N}}=256.5$ K 的长程三维
反铁磁有序，但远高于此温度时自旋涨落由二维涨落主导。三维有序在非弹性中子散射
数据中给出尖锐的布拉格峰，而关联长度 $\xi$ 的短程有序给出宽度正比于 $\xi^{-1}$
的宽峰（长程有序时 $\xi\to\infty$，峰变得非常尖锐）。实验数据显示峰随材料冷却而
变窄；不同入射能量的数据与基于二维量子非线性 sigma 模型预言的计算曲线一同示于
图 8.13。该理论除从拉曼散射测得的 $\mathsf{J}$ 外没有可调参数——符合程度因此
极为惊人，表明这一体系开始得到细致的理解。

\begin{figure}
\centering
\includegraphics[width=0.5\textwidth]{fig8_13.pdf}
\caption{由非弹性中子散射导出的 $\mathrm{Sr}_2\mathrm{CuO}_2\mathrm{Cl}_2$ 中
的逆关联长度 $\xi^{-1}$。数据对应不同初始中子能量 $E_i$，理论曲线基于二维量子
非线性 sigma 模型的预言。取自 Greven 等 1994。}
\end{figure}

尽管如此，把这些结果应用于高 $T_{\mathrm{c}}$ 超导电性问题仍有争议。近期诱人的
结果表明：某些体系中掺杂空穴并非随机排布，而是分凝成条纹（stripes），
条纹的排布相当精微地依赖掺杂水平。这一引人入胜的协同行为看起来是谜团的一部分——
但它是拼图的关键一块还是又一个转移视线的干扰项，仍待观察。

本节迄今处理的是维数 D=3 的海森伯自旋——尽管它们位于 d=2 的晶格上，自旋可以
指向三维空间任意方向。若磁矩被约束在平面内，我们得到二维 XY 模型（D=2）。此时
$T\neq0$ 没有长程序，但仍有一个有趣的相变可以发生：低温下自旋--自旋关联函数按
代数方式衰减，而高温下按指数衰减。两个区域之间有转变温度，与涡旋（vortex）的
热稳定性相联系。

\begin{figure}
\centering
\includegraphics[width=0.4\textwidth]{fig8_14.pdf}
\caption{XY 模型中的涡旋态。}
\end{figure}

在 XY 系统中产生一个涡旋（如图 8.14 所画）的能量近似为 $\pi\mathsf{J}\ln(R/a)$，
R 是系统尺寸、a 是晶格常数。这一能量代价在 $R\to\infty$ 时发散，看起来涡旋的
产生非常昂贵。然而涡旋中心可以位于系统 $(R/a)^{2}$ 个格点中的任何一个，故涡旋
的熵为 $S=k_{\mathrm{B}}\ln[(R/a)^{2}]$。于是涡旋的自由能为
\begin{equation*}
F = (\pi\mathsf{J} - 2k_{\mathrm{B}}T)\ln(R/a)\tag{8.3}
\end{equation*}
在温度
\begin{equation*}
T_{\mathrm{KT}} = \frac{\pi\mathsf{J}}{2k_{\mathrm{B}}},\tag{8.4}
\end{equation*}
以上变为负——这就是 Kosterlitz--Thouless 转变（Kosterlitz--Thouless
transition）：涡旋可以被热自发产生。这是一种特殊的相变，因为它发生在两个完全
无序的相之间——一个在低温、一个在高温。它伴随刚性的改变：低温态有弹性刚度而
高温态没有。这一转变因此也可以看作涡旋解缚转变（vortex unbinding
transition），因为 $T_{\mathrm{KT}}$ 以下系统中的涡旋都被强烈束缚。它被称为
拓扑相变（topological phase transition），因为没有对称性破缺（与铁磁--顺磁
转变相对照）。

若自旋被约束在一条线上，我们得到二维伊辛模型（d=2、D=1）——它确实有磁相变。
$K_2CoF_4$ 与 $Rb_2CoF_4$ 之类的材料是二维伊辛模型的近似实现，已用中子散射等
技术详细研究。

\section{量子相变}

第六章考虑的相变都由温度驱动：这类相变中，样品升温经过转变时摧毁有序的是热涨落
（其能量尺度由 $k_{\mathrm{B}}T$ 控制）。然而，若转变由其他变量控制（如压强、
磁场或掺杂水平），则在该变量的某个临界值处，可以有一个原则上能发生在绝对零度的
转变。这样的零温相变称为量子相变（quantum phase transition），其发生点称为量子
临界点（quantum critical point）。相关涨落不再是热的，而是由海森伯不确定性原理
决定的量子力学涨落。量子涨落与热涨落性质颇为不同，需要完全不同的理论来描述。
量子系统由复值波函数描述，这带来经典系统所没有的特征。

温度驱动的经典相变中，关联长度与关联时间随转变的逼近而发散：序参量涨落越来越慢、
尺度越来越大。因此存在某个与热涨落相联系的频率（设为 $\omega^{*}$），在临界
温度 $T_{\mathrm{c}}$ 处趋于零。只要临界点附近 $k_{\mathrm{B}}T_{\mathrm{c}}\gg\hbar\omega^{*}$，
临界涨落将表现得经典。因此对量子相变（$T_{\mathrm{c}}=0$），量子涨落不可忽略。
预期具有量子临界点的体系表现由量子涨落控制的反常动力学，可用 NMR 与中子散射等
技术探测。

量子相变的一个例子是伊辛磁体 $LiHoF_4$：沿垂直于伊辛自旋易轴的方向加磁场，可以
在绝对零度摧毁其铁磁有序（见图 8.15）。这一反直觉的行为之所以发生，是因为磁场
便利了上下自旋态之间的量子隧穿。超过临界磁场，量子涨落足以摧毁铁磁有序。这里的
顺磁态不同于通常的高温顺磁体——后者中自旋不停地（但经典地）在上下态之间涨落；
$LiHoF_4$ 中基态有一个唯一的相位相干波函数——上下自旋态的量子叠加。

\begin{figure}
\centering
\includegraphics[width=0.45\textwidth]{fig8_15.pdf}
\caption{$LiHoF_4$ 的相图。虚线是只计入电子自旋自由度的平均场理论；实线还纳入了
核超精细相互作用。取自 Bitko 等 1996。}
\end{figure}

有若干材料表现涉及相图中磁性与超导区域相邻的量子临界点。其中包括重费米子
\fn{2}{重费米子化合物的有效质量非常大，因而得名；它们常是 Ce 或 U 的化合物。}化合物如 CePd$_2$Si$_2$（见图 8.16(a)）——常压下是反铁磁体，施加静水压\fn{3}{静水压是各向同性施加的压强，沿所有方向相同。}后可以变为超导。20 kbar、T=0 处的相变就是一个量子临界点。

第二个例子是基于有机分子 ET\fn{4}{ET 是 BEDT-TTF 的缩写，而 BEDT-TTF 又是 bis-ethylenedithiotetrathiafulvalene（双亚乙基二硫代四硫富瓦烯）的缩写。}的一族
有机超导体。可以制备很纯的 ET 电荷转移盐晶体：一对 ET 分子共同向无机阴离子
（如 $\mathrm{Cu(NCS)_2^-}$）捐出一个电子。晶体结构由 ET 分子层与无机阴离子层
交替堆叠构成。由于 ET 分子的分子交叠，有机层面内电导率高，而层间电导低。有机
材料比无机材料软得多，静水压的效果远大于在 $CePd_2Si_2$ 中。而且，通过改变
阴离子（例如使其更臃肿），可以用化学手段施加“负压”。图 8.16(b) 表明：联合
化学压强与外加压强可以得到相似的相图——同样，反铁磁相紧邻超导相。

第三个例子已经描述过：反铁磁体 $La_2CuO_4$ 掺 Sr 后变为超导。这里控制性质的
变量是 Sr 掺杂水平 x。在临界 x 值处系统开始从反铁磁切换到超导。

相图中磁性\fn{5}{在 UGe$_2$ 中，相图的超导区紧邻铁磁区（Saxena 等 2000）。}区域与超导区域的这种邻近是如下判断的证据：在这些材料的一部分中，传递超导电性的是自旋涨落而非声子。\fn{6}{在超导电性的 BCS 理论中，电子之间的吸引配对相互作用由声子传递。这对常规超导体看来正确，但对重费米子、有机与高温超导体则不是正确的描述。}然而在写作之时，这一故事仍远未收尾。

\begin{figure}
\centering
\includegraphics[width=0.4\textwidth]{fig8_16.pdf}
\caption{（a）重费米子超导体 $\mathrm{CePd}_2\mathrm{Si}_2$ 的温度--压强相图
（取自 Mathur 等 1998）：AF=反铁磁，S=超导，M=金属。（b）有机超导体
$\kappa$-(ET)$_2$Cu(NCS)$_2$、$\kappa$-(ET)$_2$Cu[N(CN)$_2$]Br、
$\kappa$-(ET)$_2$Cu(CN)[N(CN)$_2$] 与 $\kappa$-(ET)$_2$Cu[N(CN)$_2$]Cl 的
温度--压强相图：压强轴包括化学改变单胞尺寸所致的“化学压强”效应以及常规静水压
（改编自 Kanoda 1997、Lubczynski 等 1996）。（c）铜氧化物超导体
$\mathrm{La}_{2-x}\mathrm{Sr}_x\mathrm{CuO}_4$ 的相图（改编自 Kastner 等
1998）：SG=自旋玻璃。}
\end{figure}

\section{薄膜与多层膜}

迄今考虑的是出现在晶体中的二维磁体。但也可以用表面科学技术直接生长极薄的磁性
薄膜来人为构造这类体系。用分子束外延（molecular-beam epitaxy），可以在平整
衬底上生长厚度仅单个原子层的磁性薄膜。生长在超高真空条件下进行，尽量减少污染与
氧化。衬底的晶格常数与对称性选得与被生长材料紧密匹配。不仅能生长薄膜，还能
生长多层膜：可以制备三明治结构，经由非磁隔离层的磁耦合研究将在后面一节考虑。

单层膜的磁性质预计与生长它的块体材料大不相同。膜表面的原子比块体中的原子最近邻
更少。最近邻的减少收窄电子带宽、从而提高费米能级处的态密度 $g(E_{\mathrm{F}})$，
提高满足 Stoner 判据（见 7.3 节）的机会，从而增强磁有序的倾向。因此磁性薄膜
被预期具有增强的磁矩。此外，预计 Pd 或 V（通常非磁的金属）的单层会因此变得
铁磁。

这些技术能耍的另一个把戏是：恰当选择生长薄膜的衬底，生长磁性元素的亚稳相。
在晶体衬底上生长磁性薄膜时，界面处存在的力可以把膜驱入不同的晶体学相。这一替代
相可能已经是已知的高温相或高压相，也可能是在块体中完全未知的相。室温下 Ni 通常
是面心立方（fcc），Fe 通常体心立方（bcc），Co 通常六角密堆（hcp）。然而用彼此
晶格匹配的特定衬底，可以生长亚稳晶体学相：例如亚稳的 fcc Co 可以稳定在 fcc Cu
表面，bcc Co 可以稳定在 GaAs(110) 表面。由于与晶体结构改变相联系的能量
（$\sim0.1$ eV/原子）和磁结构改变（如铁磁变反铁磁）相联系的能量同量级，薄膜的
磁性质常常急剧依赖生长条件以及衬底的结构。

薄膜表面原子对称性的改变也影响磁各向异性与易磁化方向。取最低阶，铁磁层的各向
异性能量可写为
\begin{equation*}
E_{\mathrm{an}} = K\sin^{2}\theta\tag{8.5}
\end{equation*}
K 是有效各向异性常数，为三项之和：
\begin{equation*}
K = \frac{2K_{\mathrm{s}}}{t} + K_{\mathrm{v}} - \mu_0M^{2},\tag{8.6}
\end{equation*}
$\theta$ 是磁化强度与表面法线的夹角。式 8.6 第一项代表表面各向异性：$K_{\mathrm{s}}$
是表面（或界面）各向异性常数，t 是层厚；因子 2 的出现是因为每层一般有两个面。
第二项是体积各向异性，可源于晶格应变，或因为我们有一个轴垂直于膜面的单轴单晶；
$K_{\mathrm{v}}$ 是体积各向异性常数。第三项是膜的形状各向异性（见 6.7.2 节）：
它是偶极贡献，按“磁极在平面上均匀分布”的假定计算；实践中界面粗糙度可以使其
显著低于这一计算值。

厚膜中偶极项主导，磁化强度位于膜面内。薄膜中表面（界面）项可能主导（因为
$t^{-1}$ 因子），自发磁化强度可能变得垂直于膜面。这一垂直各向异性（perpendicular
anisotropy）在记录应用中特别有用。\fn{7}{记录应用中，膜上的小区域存储信息比特。这些小区域因邻居而感受的偶极场，在磁化垂直于膜面时较低，因此比特原则上可以以更高的密度存储。不过这一解释略过了大量实际问题——它们有时会使制造商偏爱磁化位于面内的体系。}$K_{\mathrm{s}}$ 的来源是表面态对称性更低——增强了表面的磁晶各向异性。它受界面
粗糙度以及每一界面处晶格失配的影响。制备具有大垂直磁化强度的体系对磁记录技术的
发展极为重要。

已有多种实验技术用于研究薄膜的磁化强度。穆斯堡尔技术已在 3.2.3 节描述；
磁光方法将在下一节介绍。极化中子反射（polarized neutron reflection，PNR）
是自旋极化中子以掠入射从磁性多层膜上反射的技术：自旋向上与向下中子的反射率
不同，可用来推断随深度变化的磁化强度（方向与大小）。磁力显微术（MFM）中，
悬臂上的磁矩扫过样品表面：只要样品上方磁场存在梯度，针尖就会受力，因此该技术对
畴壁敏感（见 6.7.6 节）。另一种显微技术是 SEMPA（扫描电子显微术与极化分析）：
常规扫描电子显微镜研究二次电子的发射（二次电子保持其自旋极化，给出表面磁化强度
的信息）。其他技术包括克尔显微镜（偏振激光束扫过表面以对磁化强度成像）与扫描
SQUID 显微术。

\section{磁光学}

磁光效应最早由 Michael Faraday 于 1845 年研究：他证明偏振光通过置于磁场中的
玻璃块后出射时偏振面被旋转。这一效应如今称为法拉第效应（Faraday effect）。
相关现象由 John Kerr 于 1877 年发现：法拉第效应发生于透射，磁光克尔效应
（magneto-optic Kerr effect）发生于反射。John Kerr 把光从电磁铁抛光的磁极上
反射，注意到偏振面被旋转。

两种效应都与自旋--轨道相互作用相关，严格的处理需要微扰论。不过一般原理可以这样
相当简单地理解：右旋与左旋圆偏振光使材料中的电荷向相反方向旋转，因此每种偏振
对轨道角动量的贡献符号相反。磁场沿磁场方向产生自旋极化，自旋--轨道相互作用于是
给两种圆偏振带来大小相同、符号相反的能量贡献。这导致右旋与左旋偏振在材料中折射率
不同。若平面偏振光入射磁性材料，应把它看作右旋与左旋圆偏振光束之和——二者以
不同速度在材料中传播；出射时两束重新组合，但它们之间的相位滞后意味着出射束的
偏振面被旋转。

各向同性介质中介电张量成为
\begin{equation*}
\epsilon = \epsilon\begin{pmatrix} 1 & \mathrm{i}Q_z & -\mathrm{i}Q_y\\ -\mathrm{i}Q_z & 1 & \mathrm{i}Q_x\\ \mathrm{i}Q_y & -\mathrm{i}Q_x & 1 \end{pmatrix},\tag{8.7}
\end{equation*}
$\mathbf{Q}=(Q_x,Q_y,Q_z)$ 称为 Voigt 矢量，沿磁场方向排列，其大小依赖材料。
介电张量中的非对角项正是磁光效应的来源。该张量给出两个圆偏振简正模，介电常数
$\epsilon_{\pm}=\epsilon(1\pm\mathbf{Q}\cdot\hat{\mathbf{k}})$，$\hat{k}$ 是光
的传播方向。圆模式以不同速度行进、以不同方式衰减。出射波组合起来给出旋转的偏振
轴以及椭圆度（不同的衰减使出射光的偏振略呈椭圆）。
\bio{Woldemar Voigt\\（1850--1919）}

金属磁体的研究用克尔效应而非法拉第效应，因为金属吸收光、不能透射研究——除非
样品非常薄。克尔效应对磁性薄膜与表面特别有用；此时该技术称为 SMOKE（表面磁光
克尔效应，surface magneto-optic Kerr effect）。它可以在超高真空（UHV）腔内的
薄膜生长过程中使用：激光束穿过腔窗入射、从样品反射、再经另一窗口出腔。光偏振的
变化可用正交偏振器监测。

克尔效应可以在多种几何下进行（见图 8.17）：极向克尔效应（polar Kerr effect）中
磁化方向 M 垂直于膜面。若 M 在膜面内，M 在光的散射面内时可测纵向克尔效应
（longitudinal Kerr effect）；M 垂直于入射面时可测横向克尔效应（transverse
Kerr effect）。由于可在生长过程中进行，SMOKE 对表面磁性、膜的磁性随膜厚的变化
以及表面各向异性的研究很有用（例如它可以用来显示磁化在什么膜厚从面内变为面外）。

磁光效应在磁记录中极为有用：可以让偏振激光束极快地扫过盘面，通过反射光偏振的
变化读出盘上编码的磁信息。虽然用磁光效应读盘，写数据却用另一束激光把盘上微小
区域局部加热超过 $T_{\mathrm{C}}$ 并对盘加磁场；只有受热区域翻转——盘上其余
信息不变。激光束然后移到下一区域写下一位。在 21 世纪之交，磁光盘已开始具有巨大
的重要性：提供了一种便宜而便携的存储较大量数据的方法。

\begin{figure}
\centering
\includegraphics[width=0.3\textwidth]{fig8_17.pdf}
\caption{（a）极向克尔效应；（b）纵向克尔效应；（c）横向克尔效应。}
\end{figure}

\section{磁电阻}

材料在磁场 H 作用下电阻 R 的变化称为磁电阻（magnetoresistance）。磁电阻
$\Delta\rho/\rho$ 通常定义为
\begin{equation*}
\frac{\Delta\rho}{\rho} = \frac{R(H)-R(0)}{R(0)}\tag{8.8}
\end{equation*}
\mnote{我们将看到，许多有趣的磁电阻是大的负效应。把磁电阻定义为 $(R(H)-R(0))/R(H)$——即除以较小的数——可以使你的新型“神奇磁性材料”的磁电阻显得大得出奇。这种把戏在研究文献中不止一次被使用过！}
这是一个技术上有用的量，因为磁电阻传感器被广泛使用（例如读取信用卡磁条的磁场）。
自由电子气不表现磁电阻——即使有效质量各向异性；磁电阻只出现在多于一种载流子的
模型中，其高场行为依赖费米面的拓扑。该效应由开尔文勋爵（Lord Kelvin，当时名
William Thomson）于 1856 年首先发现：他检验铁样品的电阻，发现纵向加磁场时铁的
电阻增加 0.2\%，横向加场时减少 0.4\%。
\bio{William Thomson（Lord Kelvin）\\（1824--1907）}

\begin{figure}
\centering
\includegraphics[width=0.45\textwidth]{fig8_18.pdf}
\caption{铁磁体中的自旋劈裂能带。}
\end{figure}

\subsection{铁磁体的磁电阻}

铁磁体中负磁电阻的观察曾是非常令人困惑的。金属承载电流时，电子向费米面不同部位
的迁移使散射最小——电子找到跨样品的耗散最小路径。因此若迫使电子走别的路径（例如
因为加了磁场），它们会走散射更多的路径。一般预期正磁电阻。低温下、厚度小于平均
自由程的样品中有时可以观察到负磁电阻：若磁场在面内且垂直于电流方向，电子路径
描绘直径更小的轨道，表面散射减少。但在铁磁体中，负磁电阻的解释必须完全不同。

Mott 对这一问题提供了非常重要的洞见：他考虑 Ni 的输运性质（自旋劈裂能带示意于
图 8.18），在 Ni 中把组态从 $(3d^84s^2)$ 变到 $(3d^94s^1)$ 或 $(3d^{10})$ 只需
几 eV。一般把 Ni 视为 $(3d^{9.4}4s^{0.6})$。d 带很窄（这是过渡金属铁磁性的必要
条件——使 $g(E_{\mathrm{F}})$ 大、Stoner 判据满足），故 $m_{\mathrm{d}}^{*}\gg m_{\mathrm{e}}$；
s 带近乎自由，$m_{\mathrm{s}}^{*}\sim m_{\mathrm{e}}$。于是电导率
\begin{equation*}
\sigma = \frac{n_{\mathrm{s}}e^{2}\tau_{\mathrm{s}}}{m_{\mathrm{s}}^{*}} + \frac{n_{\mathrm{d}}e^{2}\tau_{\mathrm{d}}}{m_{\mathrm{d}}^{*}},\tag{8.9}
\end{equation*}
由第一项主导，导电主要来自 s 电子。式 8.9 中 $n_{\mathrm{s}}$ 与 $n_{\mathrm{d}}$
分别是 s 带与 d 带中的电子数，散射时间 $\tau_{\mathrm{s}}\sim\tau_{\mathrm{d}}$。

跃迁概率主要来自 s$\to$d 跃迁。低温（$T\ll T_{\mathrm{c}}$）下所有未占据的 d 态
都是反平行的，只有一半的 s 电子能跃迁；$T>T_{\mathrm{c}}$ 时所有 s 电子都能跃迁，
散射更多。因此预期 $T_{\mathrm{c}}$ 以下电阻率下降。现在考虑对 Ni 加磁场：磁场
可能增强自旋极化、减少 s$\to$d 跃迁——因此观察到负磁电阻。

多数弹性碰撞过程中电子保持其自旋，这些碰撞由相对较短的弛豫时间 $\tau$ 刻画。
然而，即使弱的自旋--轨道耦合也允许弱的自旋翻转散射，其弛豫时间长得多，为
$\tau_{\mathrm{sf}}$。无外力时，自旋向上与向下电子平衡分布中产生的微扰先在时间
$\tau$ 内衰减成每种自旋内的均匀分布，再在时间 $\tau_{\mathrm{sf}}$ 内达到平衡
分布。$T\ll T_{\mathrm{c}}$ 时自旋翻转散射预期不占主导，铁磁体可以很好地用双电流
模型（two-current model）近似——$\uparrow$ 与 $\downarrow$ 电子电流被独立处理。
它在描述少量一种过渡金属（杂质）溶入另一种过渡金属（宿主）的合金性质上特别成功。
某些过渡金属杂质引起的散射强烈依赖自旋——这是宿主 d 带自旋劈裂、杂质 d 能级自旋
劈裂以及宿主与杂质态在自旋 $\uparrow$ 与 $\downarrow$ 方向不同杂化共同作用的
结果。例如 Fe 中的 Cr 杂质对自旋 $\uparrow$ 电子的散射强得多，导致两自旋态电阻率
之比 $\rho_{\uparrow}/\rho_{\downarrow}\sim6$。高温下与自旋波碰撞导致的自旋翻转
散射引起自旋混合（spin-mixing）——两个自旋通道之间区别的模糊。这一概念应记在
心里，它将在三明治结构巨磁电阻的讨论中回归。

\subsection{各向异性磁电阻}

铁磁体中测得的磁电阻依赖磁化强度相对材料中电流方向的取向——这一效应称为各向
异性磁电阻（anisotropic magnetoresistance）。其起源与自旋--轨道相互作用及其对
s-d 散射的影响相关。自旋--轨道相互作用降低原子波函数的对称性并混合各态；晶轴
决定 L 的方向、磁化强度决定 S 的方向，态的混合导致各向异性散射。用对称性论证，
可以预言磁电阻对磁化强度与电流密度方向依赖的一般形式（对特定晶体或多晶材料）。

\subsection{巨磁电阻}

各向异性磁电阻效应相当小，因此 1980 年代末发生了一场小小的革命：在 Fe/Cr/Fe
多层膜中发现了非常大的效应（命名为巨磁电阻，giant magnetoresistance，GMR）。
低温强场下发现了超过 50\% 的大负磁电阻（见图 8.19）。该效应与磁 Fe 层反铁磁耦合
的样品相联系。人们发现：磁性层之间经隔离层的耦合随隔离层厚度增加而振荡变号；
某些厚度下是铁磁的，厚度更大时变为反铁磁，然后又回到反铁磁。\footnote{原书如此。按耦合随隔离层厚度振荡变号的物理，此处应为“回到铁磁”——疑为原书排印之误。——译注}振荡周期看来是十个
原子间距的量级（通常在 9 到 18 \AA{} 之间），其值主要由隔离层金属决定、与铁磁
金属无关。隔离层与铁磁晶格的良好匹配有利于大的耦合。在最好的情形（如 Co/Ru
超晶格），振荡幅度随隔离层厚度 t 以 $1/t^{2}$ 衰减。

\begin{figure}
\centering
\includegraphics[width=0.5\textwidth]{fig8_19.pdf}
\caption{Fe/Cr/Fe 多层膜中的巨磁电阻（Baibich 等 1988）。}
\end{figure}

这些振荡的起源与块体金属中相距 r 的两个局域自旋之间的 RKKY 相互作用
$\mathsf{J}(r)\sim\cos(2k_{\mathrm{F}}r)/r^{3}$ 相关。对界面上所有自旋求和后，
耦合变为 $\mathsf{J}\sim\cos(2k_{\mathrm{F}}t)/t^{2}$（t 是两铁磁层的间距）。
耦合的振荡可以直接联系于隔离层费米面的拓扑。若设铁磁金属具有全满的多数自旋 d 带，
则铁磁耦合时少数自旋空穴被束缚在隔离层内，反铁磁耦合时没有束缚。粒子数守恒下，
不同耦合之间的能量差完全由隔离层（厚度 t）中少数自旋空穴能量的尺寸量子化决定：
隔离层中的能级是分立的，态密度由一系列台阶组成，每个台阶的宽度正比于 $1/t^{2}$。
t 增大时台阶变窄，总有一个台阶最终扫过费米能级。用于计算这一效应的形式理论可以
理解为德哈斯--范阿尔芬效应的一维类比：此处尺寸量子化源于垂直于层面方向上的一维
（而非二维）束缚。

三明治结构中磁耦合的类型可以直接影响观察到的磁输运行为——后者对磁层的排列极其
敏感，反铁磁耦合时 GMR 效应最大。对 GMR 的第一个解释基于双电流模型（见上），
分别考虑 $\uparrow$ 与 $\downarrow$ 电子的各自电流（$\uparrow$ 指平行于多数
自旋带）。这一讨论中我先设平均自由程 $\lambda$ 远大于 Cr 层间厚度 $t_{\mathrm{Cr}}$。
想象一个 Fe/Cr/Fe 结构：零场中两 Fe 层的磁化反平行排列，并设
$\rho_{\downarrow}\ll\rho_{\uparrow}$。有两种情形：

（1）$H>H_{\mathrm{s}}$（$H_{\mathrm{s}}$ 为饱和场）。此时 Fe 层中的磁矩平行
排列（见图 8.20(a)），电阻率 $\rho$ 为
\begin{equation*}
\frac{1}{\rho} = \frac{1}{\rho_{\uparrow}} + \frac{1}{\rho_{\downarrow}} \Rightarrow \rho\sim\rho_{\downarrow}.\tag{8.10}
\end{equation*}
散射较少的电子形成有效短路。

（2）H=0。两铁层中的磁矩现在反平行（见图 8.20(b)）。此时电子相对局域磁化在
各层中交替地为 $\uparrow$ 与 $\downarrow$，自旋 $\uparrow$ 与 $\downarrow$ 通道
被有效地“混合”（对比上文考虑的高温合金中自旋波所致的自旋混合），于是
$\rho_{\uparrow}\to\rho_{\mathrm{av}}$、$\rho_{\downarrow}\to\rho_{\mathrm{av}}$
（$\rho_{\mathrm{av}}=(\rho_{\uparrow}+\rho_{\downarrow})/2$），总电阻率 $\rho$
为
\begin{equation*}
\rho = \frac{\rho_2+\rho_1}{4} \gg \rho_1.\tag{8.11}
\end{equation*}
这再次预言了观察到的负磁电阻。巨大效应来自 $\rho_{\uparrow}$ 与 $\rho_{\downarrow}$
的不等（可以非常大），以及“自旋混合”可以简单地通过施加磁场开启与关闭的便利。

\begin{figure}
\centering
\includegraphics[width=0.6\textwidth]{fig8_20.pdf}
\caption{（a）$H>H_{\mathrm{s}}$ 时 Fe 层中的磁矩平行。（b）铬层厚度
$t_{\mathrm{Cr}}$ 选择得当时，$H=0$ 时 Fe 层中的磁矩反平行。}
\end{figure}

对更厚的 Cr 层，界面处的自旋依赖散射影响界面附近厚度约为平均自由程 $\lambda$
的层内的电子分布函数。若 $t_{\mathrm{Cr}}\gg\lambda$，预期磁电阻大致按
$\exp(-t_{\mathrm{Cr}}/\lambda)$ 下降。但若没有反铁磁耦合就不会有巨磁电阻效应；
既然耦合随 $t_{\mathrm{Cr}}$ 周期出现，磁电阻也预期随层间厚度增大而振荡，同时
指数衰减。

通过多层膜的磁耦合可以用 8.7 节描述的许多磁性薄膜实验技术测量。有用的补充是
铁磁共振（铁磁体的 NMR 或 ESR 类比）与布里渊光散射技术（用光的非弹性散射研究
激发）。二者都可用于研究磁性薄膜中的自旋波与静磁模。多层膜中的自旋经交换、偶极
与各向异性相互作用耦合在一起。作为这一磁性系统本征模的自旋波激发，其色散关系可以
相当敏感地依赖交换耦合、各向异性与磁弹效应。

制备多层膜并非实现 GMR 的唯一途径。它在非均相 Cu-Co 合金膜中也被观察到：
Cu 基体内富 Co 晶粒中磁矩的相对取向决定磁电阻，并可由外加场改变。但 GMR 不出现
在缺乏孤立的、具有合适尺寸的大型富磁晶粒的均匀合金中。合金比多层膜容易制备，因此为传感器
应用提供了激动人心的前景。

\subsection{交换各向异性}

交换各向异性（exchange anisotropy，或单向各向异性 unidirectional anisotropy）
是可以在铁磁体与反铁磁体之间观察到的界面交换。若铁磁体的居里温度 $T_{\mathrm{C}}$
高于反铁磁体的奈尔温度 $T_{\mathrm{N}}$，则把一个沉积在另一个之上（适当地选择
铁磁体与反铁磁体，且铁磁层足够厚），并在外场中冷却通过奈尔温度后，在 $T\ll T_{\mathrm{N}}$
测得的磁滞回线看起来被移动了——好像除了外加磁场之外还有另一个磁场存在。看起来
铁磁膜沿某个方向（冷却时的方向）磁化在能量上比另一方向更有利。
\begin{figure}
\centering
\includegraphics[width=0.4\textwidth]{fig8_21.pdf}
\caption{被交换各向异性移动的磁滞回线。没有交换各向异性时回线将以零磁场为中心。}
\end{figure}
有时说它被反铁磁层交换偏置（exchange biased）了。交换偏置的效果是产生一个单向
交换场 $B_{\mathrm{ex}}$，它可以与外场 B 竞争（磁层的总能量为
$-\mathbf{M}\cdot(\mathbf{B}+\mathbf{B}_{\mathrm{ex}})$）。若 B 与
$B_{\mathrm{ex}}$ 同向，效果只是移动磁滞回线（见图 8.21）；若成直角，得到难磁化
轴响应。两种情形下，每个磁场值都存在唯一使能量极小的磁化角度。

这一技术可用于给磁电阻传感头加上“交换偏置”场，把它偏置到线性区；也用于约束
三明治结构中一个软铁磁层的磁化方向，另一层的磁化则可由外场转动，从而以严格受控
的方式测量磁电阻与耦合能量。这也是自旋阀（spin valve，见图 8.22）的基础——一种
巨磁电阻传感器：两个磁层夹非磁隔离层，其中一个磁层邻接反铁磁层。三明治结构两端
的电阻于是对沿特定方向的外磁场之值敏感。

\begin{figure}
\centering
\includegraphics[width=0.42\textwidth]{fig8_22.pdf}
\caption{自旋阀。反铁磁体（AFM）把上层铁磁层（FM1）的磁矩沿特定方向“夹持”；
下层磁层（FM2）的磁矩可自由转动、可与外加磁场排齐。两铁磁层磁矩间的夹角控制器件
的电阻。零场下，层的相对排列由非磁隔离层（NM）的厚度控制。}
\end{figure}

\subsection{庞磁电阻}

锰氧化物的输运性质近来引起了极大的兴趣。LaMnO$_3$ 含有 Mn$^{3+}$ 态的 Mn——
一种扬--特勒离子；LaMnO$_3$ 表现 A 型反铁磁有序（见 5.2.4 节）。若以二价的
Sr$^{2+}$、Ca$^{2+}$ 或 Ba$^{2+}$ 离子替换 $x$ 分数的三价 La$^{3+}$ 离子，
则在 Mn 格位引入空穴。这使 $1-x$ 分数的 Mn 离子保持为 Mn$^{3+}$（3d$^4$、
t$_{2g}^{3}$e$_g^{1}$），而 $x$ 分数变为 Mn$^{4+}$（3d$^4$、t$_{2g}^{3}$e$_g^{0}$）。\footnote{原书如此。按电子计数，Mn$^{4+}$ 应为 3d$^3$（t$_{2g}^{3}$e$_g^{0}$）——疑为原书排印之误。——译注}当 $x=0.175$ 时扬--特勒畸变消失，系统变为铁磁，居里温度在室温附近；$T_{\mathrm{C}}$ 以上材料是绝缘且非磁的，$T_{\mathrm{C}}$ 以下是金属且铁磁。尤其在 $T_{\mathrm{C}}$ 附近，材料表现出极大的磁电阻效应（见图 8.23），称为庞磁电阻（colossal magnetoresistance，缩写 CMR）——“巨”这个名字已经被占用啦！CMR 的起源与金属--绝缘体转变的存在相关。

上述行为的起源部分地与双交换现象（见 4.2.5 节）相联系。Mn$^{3+}$ 离子中 t$_{2g}$
电子紧束缚于离子，而 e$_g$ 电子是巡游的。由于双交换相互作用，e$_g$ 电子在 Mn 格位
之间的跳跃只有在两个 Mn 核心自旋排齐时才被允许（实际上跳跃概率正比于
$|\cos(\theta/2)|$，$\theta$ 是两个 Mn 核心自旋的夹角）。磁场使核心自旋排齐，
因而提高电导——尤其在 $T_{\mathrm{C}}$ 附近。

实际情况其实更复杂：由于扬--特勒效应，载流子与声子相互作用。这些体系中强的
电子--声子耦合意味着：$T_{\mathrm{C}}$ 以上载流子其实是极化子（polarons）——
伴随大晶格畸变的电子。这些极化子是磁性的、在晶格中自陷。向磁态的转变可以视为
被陷极化子的解缚。电子--声子耦合还有其他标志，包括依赖磁场的结构转变与电荷有序。
电荷有序基态可以与铁磁性竞争，并在可公度掺杂值附近被增强（例如 $x=\frac{1}{2}$
处每两个 Mn 格位有一个空穴）。这些效应尚未被完全理解，仍在积极研究之中。

\begin{figure}
\centering
\includegraphics[width=0.42\textwidth]{fig8_23.pdf}
\caption{$\mathrm{La}_{1-x}\mathrm{Sr}_x\mathrm{MnO}_3$
（x=0.175）的庞磁电阻：（a）电阻率的温度依赖；（b）等温磁电阻。取自 Tokura 等
1994。}
\end{figure}

\begin{figure}
\centering
\includegraphics[width=0.42\textwidth]{fig8_24.pdf}
\caption{锰氧化物的 Ruddlesden--Popper 相。}
\end{figure}

锰钙钛矿只是表现庞磁电阻的氧化物材料之一。钙钛矿是一个称为 Ruddlesden--Popper
相的大晶体家族的成员（见图 8.24）。这些相的一般化学式为
$X_{n+1}Mn_nO_{3n+1}$（X 是镧系元素、锶或其混合），可以看作 n 层厚的钙钛矿块
\fn{8}{钙钛矿块本质上是共角 MnO$_6$ 八面体的堆叠。}的堆叠，每块之间隔着一层岩盐型的 $(\mathrm{Sr},\mathrm{Ln})_2\mathrm{O}_2$ 层——它在电与磁两方面都倾向于解耦各块。
$n=\infty$ 时实现钙钛矿化合物，而 n=1 的情形等价于高 $T_{\mathrm{c}}$ 铜氧化物
采取的 $K_2NiF_4$ 结构。锰化合物可以被制备成一系列结构（写作之时 n=1、2、3 与
$\infty$ 都已制备）。钙钛矿（$n=\infty$）中每个 MnO$_6$ 八面体被六个八面体包围；
n=2 相中这一“配位数”降为五，n=1 相中为四。最近邻数目的减少预期会造成（主要）
源于 Mn 3d 轨道的能带宽度的各向异性缩减，从而改变这些材料的电导与磁行为。因此
这些材料可以通过使用不同掺杂剂、不同掺杂水平以及调节维数而受到精细控制。

\subsection{霍尔效应}

在导电材料中沿垂直于电流方向施加磁场，会对膜中传导电子产生横向力。这一力导致
横向霍尔电压。这称为正常霍尔效应\fn{9}{通常就简单地称为霍尔效应。}——正比于外加磁场 B，因为传导电子受到的洛伦兹力是 $\mathbf{F}=e(\mathbf{E}+\mathbf{v}\times\mathbf{B})$。铁磁体中还有一个额外效应，称为反常霍尔效应
（extraordinary Hall effect），依赖磁化强度。注意它不只是“来自 M 而非 B 的
洛伦兹力”——那已包含在正常霍尔效应中，因为 $\mathbf{B}=\mu_0(\mathbf{H}+\mathbf{M})$。
经验上横向电阻率 $\rho_{\mathrm{H}}$ 为
\begin{equation*}
\rho_{\mathrm{H}} = R_0B + \mu_0R_{\mathrm{e}}M,\tag{8.12}
\end{equation*}
$R_0$ 是正常霍尔系数，$R_{\mathrm{e}}$ 是反常霍尔系数（有时称自发霍尔系数）。
正常霍尔系数 $R_0$ 相当不依赖温度，而 $R_{\mathrm{e}}$ 通常强烈依赖温度。这一
效应不仅见于铁磁体，也见于强反铁磁体或顺磁体——它只需要局域矩的存在（例如 Tb
在其铁磁态与顺磁态中都可见）。若测霍尔电阻率随场的变化，得到直线，但达到饱和场时
梯度突然改变（见图 8.25）：低场下 B 与 M 都增大，正常与反常霍尔效应都出现，梯度
由 $R_0$ 与 $R_{\mathrm{e}}$ 共同贡献；饱和场以上 M 无法再增大，梯度只来自
$R_0$。

\begin{figure}
\centering
\includegraphics[width=0.4\textwidth]{fig8_25.pdf}
\caption{铁磁体中的霍尔电阻率。}
\end{figure}

反常霍尔效应的机制与传导电子和局域矩之间的自旋--轨道相互作用相联系。载流子相对
散射中心具有轨道角动量，其符号取决于载流子从中心的哪一侧经过。由于正比于
$\mathbf{L}\cdot\mathbf{S}$ 的自旋--轨道相互作用，这一轨道角动量 L 与散射中心的
自旋角动量 S 耦合。电子于是发现从一侧经过比从另一侧经过在能量上更有利。这导致
不对称散射，从而产生横向电流——反常霍尔效应之源。

\section{有机与分子磁体}

磁性材料传统上由本质上是无机化学产物制成。然而近来基于有机化学的有机磁性材料
引起了兴趣。碳化学的丰富结构允许对每个分子作许多小调整——由此带来的可调谐性
意味着原则上可以量身定制材料以表现期望的性质。不过这只是“原则上”，因为结构--
性质关系可能证明极其复杂。凝聚态物理中研究有机材料的动力尤其来自近期三类新材料的
惊人发现：其一，导电聚合物——电导率可高达铜等常规金属，已被用于制作聚合物晶体管
与发光二极管器件；其二，超导电荷转移盐——超导转变温度可高达 15 K；其三，碳的
“足球烯”C$_{60}$——碳的一种新同素异形体，适当掺杂后可表现超导性甚至不寻常的
磁性质。

铁磁体即使在元素中也相当罕见，且全部出现在 d 或 f 区。要在有机材料上获得磁矩，
自然的做法是寻找具有未成对自旋的有机自由基。有机自由基很多，但足够稳定、能组装
成晶体结构的很少；而且即使可能，让这些自旋铁磁排列通常也不可能。因此 1990 年代
初在某些硝酮氮氧自由基（nitronyl nitroxide）分子晶体中发现铁磁性（尽管温度相当
低）特别引人瞩目。找到的第一个这类材料是对硝基苯基硝酮氮氧自由基（p-NPNN）
（见图 8.26）：仅在其一个晶相中表现直到 $T_{\mathrm{C}}\sim0.65$ K 的铁磁性。
硝酮氮氧自由基只含 C、H、N、O 元素，因此是完全有机的磁体。每个硝酮氮氧自由基
分子上有一个与两个 N--O 基团相联系的未成对自旋。分子其余部分的微小化学改变会
显著改变晶体结构，从而改变分子间交叠以及相邻分子上未成对自旋之间的磁相互作用。
因此不同化合物具有大不相同的磁基态。这些材料的转变温度大多低于 1 K。

一些在技术上最有前景的材料是分子磁体：过渡金属离子提供局域矩，有机桥连充当交换
通路。在金属离子与有机分子上都带未成对电子的材料方面也取得了进展。它们的转变
温度远高于纯有机磁体，有时高于室温。这些材料有前景的一面是光学性质：其中一些
是透明的，且变磁时变色——常与自旋转变相联系。因此它们在需要用肉眼判断磁性变化
的场合（例如电话卡或显示应用）有许多潜在应用。

\begin{figure}
\centering
\includegraphics[width=0.45\textwidth]{fig8_26.pdf}
\caption{有机铁磁体对硝基苯基硝酮氮氧自由基（p-NPNN）的分子结构。}
\end{figure}

\section{自旋电子学}

操控与放大不同自旋类型之电流的能力，是新兴的自旋电子学（spin electronics，
又称 spintronics）领域的基础——它试图用金属铁磁体、非磁金属与半导体的组合制造
器件。这一领域的基础想法都与如下事实相关：自旋向上与向下电子在铁磁体中具有不同
的迁移率。关键特征是给定铁磁体中的自旋极化程度。自旋极化 P 定义为
\begin{equation*}
P = \left|\frac{n_{\uparrow}-n_{\downarrow}}{n_{\uparrow}+n_{\downarrow}}\right|,\tag{8.13}
\end{equation*}
对元素金属铁磁体小于 1（Fe 为 0.44，Co 为 0.34，Ni 为 0.11）。但某些材料 P=1——
意味着有一个自旋劈裂带完全空着。这类材料称为半金属磁体（half-metallic magnets，
因为一种自旋的电子是金属性的、另一种自旋的是绝缘性的）；$CrO_2$、$Fe_3O_4$ 与
一些锰氧化物据信属于此类。这类材料具有很长的自旋扩散长度，可能是自旋电子学应用
中令人兴奋的材料——例如磁隧道结中用自旋极化电子隧穿过绝缘势垒。

近来出现了几种自旋晶体管的设计尝试：利用不同自旋的电子从铁磁体到半导体层的
隧穿。传统电子学只基于电子电荷、无视电子自旋。当前人们颇为乐观：在本世纪，用
磁性材料控制电子电流的自旋，可能是电子器件的有前景的新发展。

\section*{延伸阅读}

\begin{itemize}[itemsep=0.2em]
\item J.~A.~Mydosh,《Spin glasses: an experimental introduction》, Taylor and
Francis 1993 与 K.~H.~Fischer 及 J.~A.~Hertz,《Spin glasses》, CUP 1991，都包含
自旋玻璃理论与实验工作的有用论述。
\item E.~Dagotto, Reports on Progress in Physics \textbf{62}, 1525 (1999) 是
关于自旋梯的可读综述。
\item A.~M.~Tsvelik,《Quantum field theory in condensed matter physics》, CUP
1995，是包含描述低维磁性现行理论的高级教材。
\item A.~Auerbach,《Interacting electrons and quantum magnetism》,
Springer-Verlag 1994，对低维磁体的量子力学与 Haldane 能隙有生动的处理。这些专题
在 I.~Affleck, J.~Phys.: Condensed Matter, \textbf{1}, 3047 (1989) 中亦有综述。
\item M.~A.~Kastner, R.~J.~Birgeneau, G.~Shirane 与 Y.~Endoh, Rev.~Mod.~Phys.,
\textbf{70}, 897 (1998) 综述了单层铜氧化物的磁性质与高 $T_{\mathrm{C}}$ 问题。
\item E.~Manousakis, Rev.~Mod.~Phys., \textbf{63}, 1 (1991) 综述了正方晶格上
二维海森伯反铁磁体的理论。
\item P.~M.~Chaikin 与 T.~C.~Lubensky,《Principles of condensed matter
physics》, CUP 1995，是凝聚态体系统计力学的优秀信息来源。
\item 量子相变的综述见 S.~L.~Sondhi, S.~M.~Girvin, J.~P.~Carini 与 D.~Shahar,
Rev.~Mod.~Phys., \textbf{69}, 315 (1997)。
\item 侧重自旋波激发的磁性多层膜综述是 R.~E.~Camley 与 R.~L.~Stamps, J.~Phys.:
Condensed Matter, \textbf{5}, 3727 (1993)。
\item 有机磁性综述见 O.~Kahn,《Molecular magnetism》, VCH (1993)。
\item K.~De'Bell, A.~B.~MacIsaac 与 J.~P.~Whitehead, Rev.~Mod.~Phys., \textbf{72},
225 (2000) 论述了薄膜磁体中的偶极效应。
\item 磁学研究各专题的当前状况可从重要磁学会议的会议录中了解——它们不时发表在
Journal of Applied Physics、Journal of Magnetism and Magnetic Materials 与
Physica B 上。磁学研究各方面的最新观点发表在 Journal of Magnetism and Magnetic
Materials 的第 $100\times n$ 卷（n 为整数）。
\end{itemize}

\section*{本章参考文献}

\begin{itemize}[itemsep=0.2em]
\item M.~N.~Baibich, J.~M.~Broto, A.~Fert, F.~Nguyen van Dau, F.~Petroff,
P.~Etienne, G.~Creuzet, A.~Friederich 与 J.~Chazelas, Phys.~Rev.~Lett., \textbf{61},
2472 (1988)。
\item D.~Bitko, T.~F.~Rosenbaum 与 G.~Aeppli, Phys.~Rev.~Lett. \textbf{77}, 940
(1996)。
\item J.~W.~Bray, L.~V.~Interrante, I.~S.~Jacobs, J.~C.~Bonner, 载于 Extended
linear chain compounds 3, p.353, J.~S.~Miller 编, New York Plenum Press (1983)。
\item B.~R.~Coles, B.~V.~B.~Sarkissian 与 R.~H.~Taylor, Phil.~Mag. \textbf{37},
489 (1978)。
\item J.~des Cloizeaux 与 J.~J.~Pearson, Phys.~Rev. \textbf{128}, 2131 (1962)。
\item M.~Greven, R.~J.~Birgeneau, Y.~Endoh, M.~A.~Kastner, B.~Keimer, M.~Matsuda,
G.~Shirane 与 T.~R.~Thurston, Phys.~Rev.~Lett. \textbf{72}, 1096 (1994)。
\item K.~Kanoda, Hyp.~Int., \textbf{104}, 235 (1997)。
\item M.~A.~Kastner, R.~J.~Birgeneau, G.~Shirane 与 Y.~Endoh, Rev.~Mod.~Phys.
\textbf{70}, 897 (1998)。
\item B.~W.~Lovett, S.~J.~Blundell, F.~L.~Pratt, Th.~Jest\"{a}dt, W.~Hayes,
S.~Tagaki 与 M.~Kurmoo, Phys.~Rev.~B, \textbf{61}, 12241 (2000)。
\item W.~Lubczynski, S.~V.~Demishev, J.~Singleton, J.~M.~Caulfield, L.~du Croo de
Jongh, C.~J.~Kepert, S.~J.~Blundell, W.~Hayes, M.~Kurmoo 与 P.~Day, J.~Phys.:
Condens.~Matter, \textbf{8}, 6005 (1996)。
\item N.~D.~Mathur, F.~M.~Grosche, S.~R.~Julian, I.~R.~Walker, D.~M.~Freye,
R.~K.~W.~Haselwinner 与 G.~G.~Lonzarich, Nature, \textbf{394}, 39 (1998)。
\item S.~Nagata, P.~H.~Keesom 与 H.~R.~Harrison, Phys.~Rev.~B \textbf{19}, 1633
(1979)。
\item S.~S.~Saxena, P.~Agarwal, K.~Ahilan, F.~M.~Grosche, R.~K.~W.~Haselwimmer,
M.~J.~Steiner, E.~Pugh, I.~R.~Walker, S.~R.~Julian, P.~Monthoux, G.~G.~Lonzarich,
A.~Huxley, I.~Sheikin, D.~Braithwaite 与 J.~Flouquet, Nature, \textbf{406}, 587
(2000)。
\item D.~A.~Tennant, T.~G.~Perring, R.~A.~Cowley 与 S.~E.~Nagler, Phys.~Rev.~Lett.
\textbf{70}, 4003 (1993)。
\item Y.~Tokura, A.~Urushibara, Y.~Moritomo, T.~Arima, A.~Asamitsu, G.~Kido 与
N.~Furukawa, J.~Phys.~Soc.~Jpn., \textbf{63}, 3391 (1994)。
\end{itemize}
